% !TEX root = ./gzuthesis-example.tex

% 论文基本信息配置

\gzusetup{
  %******************************
  % 注意：
  %   1. 配置里面不要出现空行
  %   2. 不需要的配置信息可以删除
  %   3. 建议先阅读文档中所有关于选项的说明
  %******************************
  %
  % 输出格式
  %   选择打印版（print）或用于提交的电子版（electronic），前者会插入空白页以便直接双面打印
  %
  output = electronic,
  %
  % 标题
  %   可使用“\\”命令手动控制换行
  %
  title  = {贵州大学学位论文 \LaTeX{} 模板},
  title* = {An Introduction to \LaTeX{} Thesis Template of \\ Guizhou
            University v\version},
  %
  % 学科门类
  %   1. 学术型
  %      - 中文
  %        需注明所属的学科门类，例如：
  %        哲学、经济学、法学、教育学、文学、历史学、理学、工学、农学、医学、
  %        军事学、管理学、艺术学
  %      - 英文
  %        博士：Doctor of Philosophy
  %        硕士：
  %          哲学、文学、历史学、法学、教育学、艺术学门类，公共管理学科
  %          填写“Master of Arts“，其它填写“Master of Science”
  %   2. 专业型
  %      直接填写专业学位的名称，例如：
  %      教育博士、工程硕士等
  %      Doctor of Education, Master of Engineering
  %
  degree-category  = {工学硕士},
  degree-category* = {Master of Science},
  %
  % 培养单位
  %   填写所属院系的全名
  %
  department = {计算机科学与技术系},
  %
  % 学科
  %   1. 研究生学术型学位，获得一级学科授权的学科填写一级学科名称，其他填写二级学科名称
  %   2. 本科生填写专业名称，第二学位论文需标注“（第二学位）”
  %
  discipline  = {电子信息},
  discipline* = {Electronic Information},
  research-direction = {大数据技术与工程},
  research-direction* = {Big Data Technology and Engineering},
  %
  % 专业领域
  %   1. 设置专业领域的专业学位类别，填写相应专业领域名称
  %   2. 2019 级及之前工程硕士学位论文，在 `engineering-field` 填写相应工程领域名称
  %   3. 其他专业学位类别的学位论文无需此信息
  %
  % professional-field  = {计算机技术},
  % professional-field* = {Computer Technology},
  %
  % 姓名
  %
  author     = {王小华},
  author*    = {Wang Xiaohua},
  student-id = {2024025084},
  %
  % 学号
  % 仅当书写开题报告时需要（同时设置 `thesis-type = proposal'）
  %
  % student-id = {2000310000},
  %
  % 指导教师
  %   中文姓名和职称之间以英文逗号“,”分开，下同
  %
  supervisor  = {张安思, 副教授},
  supervisor* = {Professor Zhang Ansi},
  %
  % 副指导教师
  %
%   associate-supervisor  = {陈文光, 教授},
%   associate-supervisor* = {Professor Chen Wenguang},
  %
  % 联合指导教师
  %
  % co-supervisor  = {某某某, 教授},
  % co-supervisor* = {Professor Mou Moumou},
  %
  % 日期
  %   使用 ISO 格式；默认为当前时间
  %
  date            = {2027-06-01},
  graduation-year = {2027},
  %
  % 是否在中文封面后的空白页生成书脊（默认 false）
  %
  include-spine = false,
  %
  % 密级和年限
  %   秘密, 机密, 绝密
  %
  secret-level = {公开},
  % secret-year  = {10},
  %
  % 封面
  clc                = {TP999},
  udc                = {2024025084},
  school-code        = {10657},
}

% 载入所需的宏包

% 定理类环境宏包
\usepackage{amsthm}
% 也可以使用 ntheorem
% \usepackage[amsmath,thmmarks,hyperref]{ntheorem}

\gzusetup{
  %
  % 数学字体
  math-style = ISO,  % GB | ISO | TeX
  math-font  = xits  % stix | xits | libertinus
}

% 可以使用 nomencl 生成符号和缩略语说明
% \usepackage{nomencl}
% \makenomenclature

% 表格加脚注
\usepackage{threeparttable}

% 表格中支持跨行
% \usepackage{multirow}

% 固定宽度的表格。
% \usepackage{tabularx}

% 跨页表格
\usepackage{longtable}

% 算法
\usepackage{algorithm}
\usepackage{algorithmic}

% 量和单位
\usepackage{siunitx}

% 参考文献使用 BibTeX + natbib 宏包
% 顺序编码制
% \usepackage[sort]{natbib}
% \bibliographystyle{gzuthesis-numeric}
% 使用 BibTeX + natbib
% \usepackage[sort]{natbib}
% \bibliographystyle{plainnat}

% GB/T 7714—2025 参考文献格式
% 顺序编码制
\usepackage[sort]{gbt7714}

% 贵州大学要求参考文献表列出全部作者
\gbtsetup{
  maxbibnames = 0,
  minbibnames = 0,
  inline      = true,
}

\bibliographystyle{gbt7714-2025-numeric}


% 著者-出版年制
% \usepackage{natbib}
% \bibliographystyle{gzuthesis-author-year}


% 本科生参考文献的著录格式
% \usepackage[sort]{natbib}
% \bibliographystyle{gzuthesis-bachelor}

% 参考文献使用 BibLaTeX 宏包
% \usepackage[style=gzuthesis-numeric]{biblatex}
% \usepackage[style=gzuthesis-author-year]{biblatex}
% \usepackage[style=gzuthesis-bachelor]{biblatex}
% \usepackage[style=apa]{biblatex}
% \usepackage[style=mla-new]{biblatex}

% 声明 BibLaTeX 的数据库
% \addbibresource{ref/refs.bib}

% 定义所有的图片文件在 figures 子目录下
\graphicspath{{figures/}}

% 数学命令
% \newcommand\dif{\mathop{}\!\mathrm{d}}
\makeatletter
\newcommand\dif{%
  \mathop{}\!%
  \ifgzu@math@style@TeX
    d%
  \else
    \mathrm{d}%
  \fi
}
\makeatother

% hyperref 宏包在最后调用
\usepackage{hyperref}